%------------------------------------------------------------------------------
\section{Methodology}
\label{sec:method}

\textbf{Pipeline overview.}
The pipeline has four custom kernels: grayscale, energy, cost DP, and seam
removal (Fig.~\ref{fig:pipeline}).
Grayscale conversion uses BT.601 luma coefficients,
$Y = 0.299R + 0.587G + 0.114B$.
The energy map uses the $L_1$ gradient magnitude
\begin{equation}
e_{i,j} = \bigl|\partial I/\partial x\bigr| + \bigl|\partial I/\partial y\bigr|,
\end{equation}
computed with forward differences and border clamping.
The seam is extracted by back-tracking the \texttt{bp} array from
$j^{*} = \arg\min_{j}\,M_{H-1,j}$.
Profiling on a 1428$\times$968 image confirms the DP consumes $\bm{\sim}$\textbf{93\%}
of per-seam time; all optimization therefore targets that kernel.

\begin{figure}[t]
\centering
\begin{tikzpicture}[font=\footnotesize,node distance=3mm and 4mm]
  \tikzstyle{k}=[draw,rounded corners,minimum height=6mm,inner sep=2pt,align=center]
  \node[k] (g)  {grayscale};
  \node[k,right=of g] (e) {energy};
  \node[k,right=of e,fill=blue!12,very thick] (d) {DP cost\\\textbf{$\sim$93\%}};
  \node[k,right=of d] (r) {remove\\seam};
  \draw[-{Latex}] (g)--(e); \draw[-{Latex}] (e)--(d); \draw[-{Latex}] (d)--(r);
  \draw[-{Latex}] (r.south) .. controls +(0,-6mm) and +(0,-6mm) .. (g.south)
    node[midway,below,font=\scriptsize] {repeat $k$ times (one seam each)};
  \begin{scope}[shift={(0,-2.2)}]
    \foreach \x in {0,1,2} \node[draw,minimum size=4.2mm] (a\x) at (\x*0.5,0.5) {};
    \node[draw,minimum size=4.2mm,fill=blue!12] (b1) at (0.5,0) {};
    \draw[-{Latex}] (a0.south) -- (b1.north west);
    \draw[-{Latex}] (a1.south) -- (b1.north);
    \draw[-{Latex}] (a2.south) -- (b1.north east);
    \node[right=5mm of a2,align=left] {row $i\!-\!1$ (in shared mem)\\$M_{i,j}\!=\!e_{i,j}\!+\!\min(\cdot,\cdot,\cdot)$\\row $i$ depends on 3 cells above};
  \end{scope}
\end{tikzpicture}
\caption{The four-kernel pipeline (top) and the DP's 3-way row dependency
(bottom). The $H$ rows are serial; the $W$ columns within each row are parallel.}
\label{fig:pipeline}
\end{figure}

\begin{figure}[t]
\centering
\includegraphics[width=\columnwidth]{figures/fig_kerneldesign.jpeg}
\caption{Kernel paradigm comparison.
\textbf{Naive} launches $H$ kernels (one per row) with a global-memory
round-trip each time.
\textbf{Fused} collapses all rows into one kernel: rows exchange via a
shared-memory ping-pong with \texttt{\_\_syncthreads} between them.
\textbf{Tiled} runs $K$ blocks in parallel (one strip per SM); each tile of
$T$ rows stays entirely in shared memory and only exchanges one cost row
per strip through global memory at the tile boundary, reducing kernel
launches from $H$ to $\lceil H/T\rceil$.}
\label{fig:kerneldesign}
\end{figure}

\textbf{Naive} --- one kernel launch per row: $H{-}1$ launches per seam,
with a \textbf{full global-memory round-trip} per row
(Fig.~\ref{fig:kerneldesign}, left).
Each row reads $W$ energy and $W$ cost values from DRAM and writes $W$
cost values plus $W$ back-pointers back; total per-seam traffic is
$\bm{4HW}$ \textbf{memory transactions}, all of which hit DRAM because
the cost array (one seam at 4K: ${\sim}$32\,MB) far exceeds the L2 cache.

\textbf{Fused} --- one thread block walks all $H$ rows top-to-bottom, keeping
previous and current cost rows in two \textbf{ping-pong shared-memory buffers}
of $W$ \texttt{int16} values each (Fig.~\ref{fig:kerneldesign}, centre).
One \texttt{\_\_syncthreads()} separates rows, ensuring the prior row's
costs are visible before the next row begins.
Shared-memory footprint: $2W \times 2\,\text{B}$; at 4K ($W{=}3840$)
this is $15.4$\,KB, well within the 96\,KB opt-in limit.
One launch replaces $H$, eliminating $(H{-}1)$ kernel overheads.

\textbf{GridSync} --- replacing the block barrier with a grid-wide
\texttt{grid.sync()} after each row lets multiple SMs collaborate, but each of
the $H$ device-wide barriers costs ${\sim}10\,\mu\text{s}$ against only
${\sim}1\,\mu\text{s}$ of per-row arithmetic at $W{=}1428$: \textbf{barrier-bound
at $\bm{\sim}$6\% occupancy}, so we abandon it. This negative result motivates
\emph{Tiled}'s barrier-free design.

\textbf{Prefetch} --- Nsight Compute identifies \emph{Fused}'s dominant stall as
\texttt{long\_scoreboard}: every warp waits ${\sim}$200 cycles for the global
load of the current energy row $e_{i,\cdot}$.
The fix is software pipelining: while computing $M_{i,\cdot}$ (which only
reads shared-memory cost data), each thread \textbf{issues the global load
for $e_{i+1,\cdot}$} into a register array \texttt{e\_next[]}.
By the time \texttt{\_\_syncthreads()} completes and row $i{+}1$ begins,
the load has resolved---arithmetic and memory transfer overlap:
\begin{equation}
\underbrace{\text{compute}\;M_{i,\cdot}}_{\text{shmem reads}} \;\big\|\;
\underbrace{\text{load}\;e_{i+1,\cdot}}_{\text{global read}}.
\end{equation}

\textbf{BytePtr} --- the back-pointer at each cell is one of $\{-1,0,+1\}$,
yet \emph{Fused} stores it as a 4-byte \texttt{int}.
Switching to \texttt{signed char} (1\,byte) \textbf{cuts back-store traffic by
$\bm{4\times}$}: at 4K ($W{=}3840$, $H{=}2160$), the back-pointer array
shrinks from 32\,MB to 8\,MB per seam pass.
Output is bit-identical; the seam extraction reads the same values.

\textbf{Templated} --- the $W\le 2048$ cap in \emph{Prefetch}/\emph{BytePtr}
arises because the register array \texttt{e\_next[]} must have a fixed
compile-time size.
Making columns-per-thread (\texttt{CPT} $= \lceil W/\text{blockDim}\rceil$) a
\textbf{compile-time template parameter} lets the host select the right
instantiation at launch; all prefetch arrays remain register-resident.
At 4K: $\text{CPT}=4$, registers per thread ${\approx}24$.
At ${\ge}6$K: $\text{CPT}=8$, register count exceeds the per-block budget
and the compiler spills to local memory---\textbf{Templated fails at $\ge$6K}.

\textbf{GridStride} --- \emph{Templated}'s register overflow at 8K is caused by
the fixed-size prefetch arrays.
\emph{GridStride} replaces the CPT register array with a \textbf{scalar
prefetch variable}: each thread strides over columns $t,\;t{+}B,\;t{+}2B,\ldots$
(where $B{=}\text{blockDim.x}$), issuing one scalar prefetch load per stride.
Register count stays constant at any $W$, enabling \textbf{arbitrary-width support},
at the cost of losing the multi-element register-resident overlap that
\emph{Templated} achieves---hence \emph{GridStride} shows a
20--40\% latency deficit vs.\ \emph{Templated} on images where both run.

\textbf{Tiled} --- \emph{Templated} and \emph{GridStride} both use a
\emph{single block}, leaving 79 of 80 SMs idle.
\emph{Tiled} breaks this with a \textbf{halo argument}
(Figs.~\ref{fig:kerneldesign} right, \ref{fig:tiledhalo}): the $W$ columns
are divided into $K$ strips of center width $S = \lceil W/K \rceil$.
Strip $k$ covers columns $[kS - T,\;(k{+}1)S + T)$---i.e., $T$ halo columns on
each side---so $K$ blocks run concurrently, one per SM.

The correctness argument is a seam-drift bound: within a tile of $T$ rows, the
seam shifts by at most one column per row, so after $T$ rows it has drifted by
at most $T$ columns.
Any seam error is therefore contained within the $T$-column halo, and the
center $S$ columns are \textbf{bit-exact}.
Only the center strip of the cost row is written to the global ping-pong
buffer at each tile boundary; the halo is recomputed by the neighboring block
from its own center result.

Shared memory per block stores two ping-pong cost rows and one back-pointer row,
each of width $S{+}2T$:
\begin{equation}
\text{shmem} = 2(S{+}2T)\cdot 2\,\text{B}
             + (S{+}2T)\cdot 1\,\text{B}
             = 5(S{+}2T)\,\text{B}.
\label{eq:shmem}
\end{equation}
At $K{=}60$, $T{=}64$, $W{=}3840$: $S{=}64$ and
$\text{shmem} = 5{\times}192 = 960\,\text{B}$ per block---far below the
96\,KB limit.
\textbf{Kernel launches drop from $H$ to $\lceil H/T\rceil$}
(68 at 8K with $T{=}64$) while all $K{\approx}60$ SMs stay active.
Output is bit-exact against \emph{Templated}.
A 3$\times$3 parameter sweep ($K{\in}\{40,60,80\}$, $T{\in}\{32,64,128\}$)
shows all nine combinations land within ${\sim}5\%$ of the optimum
(Appendix~\ref{app:sweep}):
\textbf{the design is robust to parameter choice}.

\begin{figure}[t]
\centering
\includegraphics[width=\columnwidth]{figures/fig_tiledhalo.png}
\caption{\textbf{Tiled} strip and halo layout.
\textit{Top}: The $W$ columns are divided into $K$ strips; each strip is
extended by $T$ halo columns on each side (orange). Adjacent strips share
halo columns without a barrier.
\textit{Bottom}: Within a tile of $T$ rows, one block computes entirely in
shared memory. The seam drifts at most one column per row, so after $T$ rows
the error stays inside the orange halo; the blue center is bit-exact.
Only the center strip is written to \texttt{d\_next} at the tile boundary.}
\label{fig:tiledhalo}
\end{figure}
